\documentclass[10pt,a4paper]{amsart}
\usepackage[margin=19mm]{geometry}
\usepackage{amsmath,amssymb,mathtools}
\usepackage[scr=rsfs,cal=euler]{mathalfa}
\usepackage{xfrac}
\usepackage[unicode,hidelinks,bookmarks=false]{hyperref}
\newcommand{\ie}{i.e.\ }
\newcommand{\cf}{cf.\ }
\newcommand{\resp}{respectively\ }
\newcommand{\Subsection}{\subsection}
\providecommand{\nobreakdash}{\nobreak\hbox{-}}
\setlength{\emergencystretch}{2em}
\multlinegap0pt
\usepackage{amssymb}
\usepackage{mathtools}
\usepackage{xcolor}
\usepackage{tensor}
\usepackage[normalem]{ulem}
\usepackage{xfrac}
\usepackage{cancel}
\usepackage{tikz}
\usepackage{braket}
\usepackage[shortlabels]{enumitem}
\newtheorem{lemma}{Lemma}
\title{Convergence of Peter--Weyl Truncations of Compact Quantum Groups}
\author{Malte Leimbach}
\date{Source: arXiv:2409.16698v2 (CC BY 4.0)}
\begin{document}
\maketitle

\noindent\emph{Source and licence.} Convergence of Peter--Weyl Truncations of Compact Quantum Groups. Version arXiv:2409.16698v2 (CC BY 4.0), distributed by arXiv under the Creative Commons Attribution 4.0 International licence, http://creativecommons.org/licenses/by/4.0/, as recorded in the arXiv licence metadata for this exact version and shown on the abstract page https://arxiv.org/abs/2409.16698v2. Full licence notice: \texttt{licenses/arxiv-2409.16698-CC-BY-4.0.txt}.\par\smallskip

\noindent\textit{Editorial scope.} Counted unit: the article's lemma lem:Equivalence-induced-operator-seminorm (source0.tex lines 1646--1653). The article's complete original proof of it is delivered with it (source0.tex lines 1655--1700, one \texttt{proof} environment of 46 lines). No mathematics is written, repaired or re-derived: the statement and the proof are the article's own lines, in the article's own notation, and no retained line is substituted or re-flowed.  No further source-local context is needed: every cross-reference of every retained line resolves inside the two delivered spans.  Nothing was omitted from any retained span, so the package carries no omission record.  No retained line contains a citation, so the wrapper carries no bibliography and no bibliography entry.  No retained line contains a figure environment, an image-inclusion command or drawing code, so no image is delivered and the package declares no asset dependency.  Editorial formatting only, in addition: an AMS \texttt{amsart} preamble that restates the article's own declarations and macros used by the retained lines, namely the environments \texttt{lemma} on separate counters, together with the standard packages those lines need.  The retained environments print in this document as Lemma 1 in order of appearance, whereas the article prints the same items with the numbering of its own full document.  The complete native source of the article as distributed in the arXiv e-print for this exact version is bundled unchanged as \texttt{sources/165782/source0.tex}.
	\begin{lemma}\label{lem:Equivalence-induced-operator-seminorm}
		For all $x \in P\mathrm{C}(G)P$, the following holds:
		\begin{align}\label{eqn:equivalence-induced-and-left-seminorm}
			\frac{1}{2}\lVert x\rVert_{\lambda,\rho}
			\leq \mathrm{Lip}^{\alpha^\tau,\beta^\tau}_{P\mathrm{C}(G)P}(x)
			\leq \lVert x\rVert_{\lambda,\rho}.
		\end{align}
	\end{lemma}

	\begin{proof}
		Let $x \in P\mathrm{C}(G)P$.
		We show that
		\begin{align}\label{eqn:Lip-rho-alpha}
			\frac{1}{2} \lVert x \rVert_\rho 
			\leq \mathrm{Lip}^{\alpha^\tau}_{P\mathrm{C}(G)P}(x)
			\leq \lVert x \rVert_\rho,
		\end{align}
		where $\mathrm{Lip}^{\alpha^\tau}_{P\mathrm{C}(G)P}(x) := \sup_{\phi \in \mathcal{S}(P\mathrm{C}(G)P)} \mathrm{Lip}_d (\phi(\alpha^\tau_\bullet(x)))$.
		To this end, using the fact that $\mathrm{Lip}_d(f) = \lVert f \rVert_\lambda$, for all $f \in \mathrm{C}(G)$, we have:
		\begin{align}
			\mathrm{Lip}^{\alpha^\tau}_{P\mathrm{C}(G)P}(x)
			&= \sup_{\phi \in \mathcal{S}(P\mathrm{C}(G)P)} \lVert \phi(\alpha^\tau_\bullet(x)) \rVert_\lambda \\
			&= \sup_{\phi \in \mathcal{S}(P\mathrm{C}(G)P)} \sup_{g \in G \setminus \{e\}} \frac{\lVert \lambda_g(\phi(\alpha^\tau_\bullet(x))) - \phi(\alpha^\tau_\bullet(x)) \rVert_\infty}{d(g,e)} \\
			&= \sup_{\phi \in \mathcal{S}(P\mathrm{C}(G)P)} \sup_{g \in G \setminus \{e\}} \sup_{h \in G} \frac{\lvert \phi(\alpha^\tau_{g^{-1}h}(x)) - \phi(\alpha^\tau_h(x)) \rvert}{d(g,e)}
		\end{align}
		By the Kadison function representation, this last quantity is bounded below and above by respectively $\frac{1}{2}$ and $1$ times 
		\begin{align}
			\sup_{g \in G \setminus \{e\}} \sup_{h \in G} \frac{\lVert \alpha^\tau_{g^{-1}h}(x) - \alpha^\tau_h(x)\rVert}{d(g,e)}
			= \lVert \alpha^\tau_\bullet(x) \rVert_\lambda.
		\end{align}
		Note that, for all $g, h \in G$, $x \in P\mathrm{C}(G)P$ and $f \in \mathrm{C}(G)$ with $\tau(f) = x$, we have
		\begin{align}
			\alpha_{g^{-1}h}(x)
			&= ((\tau \otimes \mathbf{I}^{\mathrm{C}(G)})\Delta(f))(g^{-1}h) \\
			&= \tau(p \mapsto f(pg^{-1}h)) \\
			&= \tau(p \mapsto \rho_{g^{-1}h}(f)(p)) \\
			&= \rho_{g^{-1}h}(x).
		\end{align}
		This implies that 
		\begin{align}
			\lVert \alpha^\tau_\bullet(x) \rVert_\lambda
			= \sup_{g \in G \setminus \{e\}} \sup_{h \in G} \frac{\lVert \rho_{g^{-1}h}(x) - \rho_h(x)\rVert}{d(g,e)}
			= \|x\|_\rho,
		\end{align}
		by invariance of the metric $d$.
		Altogether we obtain (\ref{eqn:Lip-rho-alpha}).
		
		Similarly, we can show 
		\begin{align}
			\frac{1}{2} \lVert x \rVert_\lambda 
			\leq \mathrm{Lip}^{\beta^\tau}_{P\mathrm{C}(G)P}(x)
			\leq \lVert x \rVert_\lambda.
		\end{align}
		Together with (\ref{eqn:Lip-rho-alpha}) this yields the claim.
	\end{proof}

\end{document}
